%======================================================================
%  Critical literature review and research gap
%======================================================================

\section{Related work and the research gap}
\label{sec:related-work}

\subsection{Corpus and source roles}

The supplied file
\texttt{Lessons from HPL and HPL-MxP on Aurora.pdf} now contains the
intended Goto et al.\ paper, \emph{Sustaining Exascale Performance:
Lessons from HPL and HPL-MxP on Aurora}, arXiv:2604.09517v1
~\cite{goto2026aurorahpl}.  Its title, authors and deployment results
were checked against the local twelve-page copy.

The four sources answer different parts of the argument.  Na et
al.\ provide the closest LLM workload precedent
~\cite{na2024llmcpu}; Ibeid et al.\ isolate memory bandwidth, latency,
mode and locality effects~\cite{ibeid2025aurora}; Goto et al.\ show the
importance of explicit tier staging and resource mapping at production
scale
~\cite{goto2026aurorahpl}; and Dongarra et al.\ provide architectural
motivation rather than a directly comparable baseline
~\cite{dongarra2026gpus}.  Their transfer limits are addressed with the
results below rather than treating the four works as equivalent
evidence.

\subsection{Na et al.: the closest inference precedent}
\label{sec:na-review}

Na et al.\ evaluate Intel Extension for PyTorch (IPEX) on a dual-socket
Xeon Max 9468 server and compare it with a previous-generation Xeon and
with single A100 and H100 GPUs running FlexGen
~\cite{na2024llmcpu}.  Their main comparisons bind one 48-core Max
socket; only the core-count sweep extends to 96 cores.  The CPU
experiments use BF16 OPT and LLaMA-2 models, prompt length 128,
generation length 32 and batch sizes from 1 to 32; a separate sweep
increases the prompt length.  This is the closest platform and workload
match in the corpus, but its results need careful qualification.

\begin{table}[H]
\centering
\small
\caption{Results from Na et al.\ and the interpretation used here
~\cite{na2024llmcpu}.}
\label{tab:na-findings}
\begin{tabularx}{\textwidth}{@{}YYY@{}}
\toprule
\textbf{Reported result} & \textbf{What it supports}
 & \textbf{What it does not establish} \\
\midrule
Depending on batch, Xeon Max geometric-mean throughput over the tested
models is 3.2--6.3$\times$ that of the Ice Lake system.
 & Matrix acceleration and the newer memory system make recent CPUs
   substantially more credible for inference.
 & A causal HBM or AMX speedup: cores, cache, DDR generation,
   microarchitecture and software paths also change. \\
\addlinespace[2pt]
Quadrant+Flat gives the best average of the four configurations under
the paper's HBM-first Flat placement.
 & In this setup, explicit HBM placement outperforms hardware-managed
   caching; unmanaged fine-grained NUMA is costly.
 & A universal preference for Quadrant mode.  Exact allocation,
   first-touch and binding details are not disclosed. \\
\addlinespace[2pt]
On average over the tested models and batches, 48 cores outperform a
naive 96-core execution.
 & Remote traffic can erase the value of additional cores.
 & An inherent one-socket scaling limit; replicas or topology-aware
   sharding are not evaluated. \\
\addlinespace[2pt]
At batch one, the CPU outperforms the tested FlexGen A100/H100 runs for
oversized models.
 & Capacity and data movement can dominate accelerator arithmetic at
   low batch.
 & General CPU superiority.  For LLaMA2-70B at batch 16 and prompt
   length at least 256, the offloaded H100 regains lower latency; when a
   model fits GPU memory, the GPUs are faster.  Multi-GPU serving is not
   tested. \\
\bottomrule
\end{tabularx}
\end{table}

The paper's GPU result is therefore best read as a data-movement result.
It compares direct CPU execution with a particular single-GPU
offloading system.  It does not compare like-for-like production
serving stacks.  The study also omits quantized weights, long
generations, energy, online arrivals, tail latency and output-quality
validation.  Repetition count, dispersion, complete software versions,
firmware, warm-up and the exact NUMA command are not reported.  Its
experimental design can be reimplemented, but the paper does not
support exact numerical reproduction.  It also proposes, without
implementing or evaluating, NUMA-aware hot/cold placement and
CPU--GPU layer partitioning.  A later placement contribution must
acknowledge that prior proposal.

\subsection{Ibeid et al.: bandwidth, latency and locality}
\label{sec:ibeid-review}

Ibeid et al.\ separate memory-system effects more clearly
~\cite{ibeid2025aurora}.  On the Borealis testbed they use Xeon Max
9470C processors and run STREAM, Intel Memory Latency Checker, PCIe and
MPI microbenchmarks, followed by three HPC applications chosen to stress
bandwidth, latency or capacity.  At the time of the study, production
Aurora was configured in Quadrant+Flat mode; the full single-socket
Flat/Cache and Quadrant/SNC4 matrix comes from Borealis.

\begin{table}[tbp]
\centering
\small
\caption{Single-socket STREAM Triad bandwidth on the Borealis Xeon Max
9470C testbed, in GB/s~\cite{ibeid2025aurora}.  The paper does not fully
specify array size, first touch or cache-state controls.}
\label{tab:ibeid-stream}
\begin{tabular}{@{}lccc@{}}
\toprule
 & \multicolumn{2}{c}{\textbf{Flat mode}} & \textbf{Cache mode} \\
\cmidrule(lr){2-3}
\textbf{Clustering} & \textbf{DDR5} & \textbf{HBM} & \textbf{DDR5+HBM cache} \\
\midrule
Quadrant, full socket & 235 & 680 & 560 \\
SNC4, full socket     & 245 & 840 & 575 \\
SNC4, one quadrant    &  60 & 210 & 150 \\
\bottomrule
\end{tabular}
\end{table}

In the two full-socket Flat-mode rows, HBM bandwidth is 2.9--3.4 times
DDR5 bandwidth, and SNC4 raises full-socket HBM bandwidth by about 24\%
over Quadrant.
However, unloaded latency favours DDR5: the paper reports 112.9\,ns for
Quadrant/DDR and 134.9\,ns for Quadrant/HBM.  Under increasing access
load, DDR5 saturates sooner while HBM sustains greater bandwidth before
latency rises sharply.  These results support a three-way distinction:
HBM for concurrent bandwidth demand, DDR5 for capacity and potentially
lower lightly loaded latency, and explicit locality for collecting the
available bandwidth.

Unless stated otherwise, Ibeid et al.\ repeat experiments three times
and report means; the latency study uses ten repetitions and reports a
mean and range.  Most other results do not report dispersion, and the
STREAM allocation and cache-state details remain insufficient for exact
reproduction.

The application results reinforce that there is no universal mode
ranking.  Explicit HBM performs best when a bandwidth-intensive working
set fits; Cache mode can help when the working set exceeds HBM; and DDR5
can remain preferable for lightly loaded latency-sensitive work.  Those
applications are not proxies for transformer inference.  They provide
mechanisms and measurements to reproduce, not a predicted LLM speedup.
Because Cache mode is direct-mapped, page-colour conflicts can also
change its result.  The paper uses page shuffling (``Fake NUMA'') to
mitigate such conflicts, so behaviour near 64\,GB must not be read as a
capacity-only effect.

\FloatBarrier

\subsection{The SNC4 result is software-dependent}
\label{sec:snc4-synthesis}

The two measurement papers appear to disagree.  Ibeid et al.\ obtain
their highest HBM bandwidth and several best application results with
SNC4 and explicit per-domain placement.  Na et al.\ find that both SNC4
configurations lose to Quadrant and report increased remote cache
traffic.  Both observations can hold: SNC4 exposes locality, but
software must align allocation and execution with that locality.

This is a productive gap rather than a contradiction in the hardware.
An MPI application can exploit SNC4 when ranks, pages and network
interfaces are explicitly bound per domain; MPI alone does not
guarantee locality.  Ibeid et al.\ recover an SNC4 MPI deficit through
explicit quadrant-to-interface mapping.  A shared-memory runtime that
allocates tensors without regard to the four domains can incur remote
traffic and smaller per-domain capacity.  CRESCO8 should therefore test
both a deliberately local SNC4 design and an unmanaged baseline if
administrators can expose that mode.

\subsection{Goto et al.: lessons from production scale}

Goto et al.\ report three production-scale Aurora campaigns, culminating in
1.012\,EF/s for FP64 HPL on 9,234 nodes, with 78.8\% parallel scaling
efficiency, and
11.64\,EF/s for mixed-precision HPL-MxP on 9,500 nodes
~\cite{goto2026aurorahpl}.  These values use different precision and are
not like-for-like throughput measurements.

The paper adds a relevant example of deliberate memory tiering.  Each
Aurora CPU has 64\,GB HBM2e and 512\,GB DDR5.  The global HPL matrix
remains in DDR because it exceeds HBM capacity, while tiles used for
panel factorisation, triangular solves and row exchange are staged into
HBM.  This demonstrates a production design that combines DDR capacity
with HBM for selected CPU-side operations.  Overall runtime is dominated
by GPU GEMM, and the paper does not ablate the staging policy, so it
does not quantify an HBM speedup transferable to LLM inference.

The study's second transferable lesson is that sustained performance
required
deterministic NUMA/PCIe-aware mapping, explicit CPU--GPU pipelining,
phase-specific communication choices and resilience to transient faults.
HPL-MxP improved from 10.6\,EF/s at ISC24 to 11.64\,EF/s at SC24,
about 9.8\%.  The authors identify AMX enablement as the primary
configuration change.  Their own impact table states that its
classifications are engineering assessments rather than controlled
ablations; the campaign comparison is therefore strong operational
evidence, not an isolated 9.8\% AMX measurement.

These findings support two methodological choices for this thesis:
resource mapping is part of the implementation, and a benchmark result
is credible only with a complete mapping and run manifest.  The paper
does not measure CPU-only LLM inference or isolate CPU HBM, so its
exascale throughput numbers should not appear as evidence for an HBM
speedup on CRESCO8.

\subsection{Dongarra et al.: architectural motivation, not a baseline}

Dongarra, Hoefler and Matsuoka argue that CPUs combining vector and
matrix engines, low-precision arithmetic and on-package HBM address the
three historical reasons for discrete GPU acceleration: arithmetic
throughput, memory bandwidth and data movement
~\cite{dongarra2026gpus}.  They retain an important boundary: GPUs
remain central to frontier-scale training, while HBM-equipped CPUs may
be attractive for inference and mixed irregular/dense scientific
workflows.

The manuscript's empirical discussion is asymmetric.  It summarizes a
companion study described as in preparation for arXiv: A64FX decode
results are measured, whereas results for a newer matrix-enhanced CPU,
including prefill and energy implications, are partly projected from
specifications.  Its trillion-parameter mixture-of-experts workload,
256K context and low-bit formats are also far from the CRESCO8 baseline.
Software maturity, energy and interconnect effects remain unresolved.
The manuscript therefore motivates the research question and the need
for runtime-aware placement, but it cannot substitute for measurements
on the Max 9480.

\subsection{What the literature establishes}

Taken together, the sources support four bounded conclusions:

\begin{enumerate}
  \item On-package HBM provides a large sustained bandwidth advantage
        over DDR5 on Xeon Max, but the achieved value depends on
        clustering and placement.
  \item Low-batch autoregressive decode is a strong candidate for
        bandwidth limitation, whereas larger prefill and batched decode
        can place more pressure on matrix throughput.
  \item Direct CPU inference can be competitive when the alternative
        repeatedly moves an oversized working set over PCIe; this is a
        capacity and data-movement claim, not a general CPU--GPU ranking.
  \item Topology results belong to the hardware--software combination.
        A poor two-socket or SNC4 result does not prove that the added
        resources are unusable.
\end{enumerate}

\subsection{What remains unanswered}
\label{sec:research-gap}

No source in the reviewed corpus performs a controlled, same-socket
comparison of local HBM and local DDR5 for modern quantized LLM
inference while separating prefill and sustained decode.  The following
questions remain unanswered by the reviewed corpus and require
measurement on CRESCO8:

\begin{itemize}
  \item how quantized weight and KV formats change the HBM benefit;
  \item what occurs when the verified live set approaches and crosses a
        64\,GB local HBM boundary;
  \item whether topology-aware sharding or independent replicas can
        reverse the reported SNC4 and two-socket regressions; and
  \item whether a performance gain also reduces energy per token under a
        clearly defined measurement boundary.
\end{itemize}

Answering the first three questions with a frozen runtime, raw data and
uncertainty is a sufficient and coherent thesis contribution.  Energy
is retained only if the measurement interface and boundary can be
validated.  GPU comparison, online serving and multi-node inference are
extensions rather than prerequisites.
